\begin{itemframe}{مسئله فروشنده دوره گرد}
\itm
ورودی مسأله‌ی فروشنده‌ی دوره‌گرد، یک گراف \textbf{کامل} و بدون جهت است که هر یال دارای یک هزینه‌ی صحیح نامنفی
$c(u, v)$
می‌باشد. هدف، یافتن یک دور همیلتونی
با کمترین هزینه است.
\itm
بیایید یک نمادگذاری دیگر به این مسئله اضافه کنیم؛ اگر A یک مجموعه از یال‌ها باشد آنگاه c(A) مجموع کل وزن یال‌های این مجموعه است.
\end{itemframe}


\begin{itemframe}{مسئله فروشنده دوره گرد}
\itm
ما به طور شهودی در بسیاری از موقعیت‌های عملی می‌دانیم کوتاه‌ترین مسیر بین دو نقطه مسیر مستقیم است. در علم ریاضیات به این مسئله نابرابری مثلثی می‌گوییم.
\itm
می‌گوییم تابع هزینه یک گراف مانند نابرابری مثلثی را ارضا می‌کند اگر
$$
c(u, w) \leq c(u, v) + c(v, w)
$$
\itm
به‌عنوان مثال، اگر رأس‌های گراف نقاطی در صفحه باشند و هزینه‌ی سفر بین دو رأس همان فاصله‌ی اقلیدسی معمولی باشد، آنگاه نابرابری مثلثی برقرار است. افزون بر این، بسیاری از توابع هزینه‌ی دیگر به جز فاصله‌ی اقلیدسی نیز نا‌برابری مثلثی را برآورده می‌سازند.
\end{itemframe}


\begin{itemframe}{مسئله فروشنده دوره گرد}
\itm
مسأله‌ی فروشنده‌ی دوره‌گرد حتی در حالتی که تابع هزینه نابرابری مثلثی را ارضا کند نیز ان‌پی-کامل است. بنابراین نباید انتظار داشته باشید الگوریتمی چندجمله‌ای برای حل دقیق این مسأله بیابید. بهتر است به دنبال یک الگوریتم‌های تقریبی باشیم.
\end{itemframe}


\begin{itemframe-s}{مسئله فروشنده دوره گرد}{حل با نابرابری مثلثی}
\itm
وزن یک درخت پوشای کمینه برای مسئله فروشنده دوره‌گرد، یک کران پایین است. بنابراین ابتدا یک درخت پوشای کمینه محاسبه می‌کنیم و این کران پایین را به‌دست می‌آوریم.
\itm
سپس از همین درخت پوشای کمینه برای ساخت یک دور همیلتونی استفاده می‌کنیم و هزینه‌ی آن حداکثر دو برابر وزن درخت پوشای کمینه خواهد بود. البته به شرطی که تابع هزینه، نابرابری مثلثی را ارضا کند.

\end{itemframe}


\begin{itemframe-s}{مسئله فروشنده دوره گرد}{حل با نابرابری مثلثی}
\itm
الگوریتم زیر این روش را پیاده‌سازی می‌کند. تابع MST-PRIM درخت پوشای کمینه را محاسبه می‌کند.
\begin{algorithm}[H]\alglr
\caption{APPROX-TSP-TOUR}
\begin{algorithmic}[1]
\Func{APPROX-TSP-TOUR}{$G, c$}
\State select a vertex $r \in V(G)$ to be a \textit{root} vertex
\State compute a minimum spanning tree $T$ for $G$ from root $r$
      using $\textsc{MST-PRIM}(G, c, r)$
\State let $H$ be a list of vertices, ordered according to when they are first visited
      in a preorder tree walk of $T$
\State \Return the Hamiltonian cycle $H$
\end{algorithmic}
\end{algorithm}
\end{itemframe}


\begin{itemframe-s}{مسئله فروشنده دوره گرد}{حل با نابرابری مثلثی}
\itm
در شبه که بالا از پیمایش پیش تریتیب
\fn{preorder tree}
استفاده شده. پیشمایش پیش ترتیب به‌طور بازگشتی هر رأس را درخت پیمایش کرده و رأس را زمانی که برای نخستین بار ملاقات می‌شود ‌(پیش از بازدید از فرزندانش) لیست می‌کند.
\itm
زمان اجرای
\texttt{APPROX-TSP-TOUR}
برابر
$\Theta(V^2)$
است.
\end{itemframe}


\begin{itemframe-s}{مسئله فروشنده دوره گرد}{حل با نابرابری مثلثی}
\itm
اکنون نشان می‌دهیم که اگر تابع هزینه‌ی یک نمونه از مسأله‌ی فروشنده‌ی دوره‌گرد نابرابری مثلثی را ارضا کند، آنگاه APPROX-TSP-TOUR دور همیلتونی‌ایی برمی‌گرداند که هزینه‌ی آن حداکثر دو برابر هزینه‌ی یک دور همیلتونی بهینه است. یعنی ضریب تقریب آن ۲ است.
\itm
\textbf{اثبات}
فرض کنید
$H^*$
یک دور همیلتونی بهینه باشد. حذف هر یال از یک دور همیلتونی، یک درخت پوشا ایجاد می‌کند و هزینه‌ی یال‌ها نامنفی است و بنابراین وزن درخت پوشای کمینه از هر دور همیلتونی‌ایی کمتر است. بنابراین اگر یک درخت پوشای کمینه را T بنامیم:
$$c(T) \leq c(H^*)$$
\end{itemframe-s}


\begin{itemframe-s}{مسئله فروشنده دوره گرد}{حل با نابرابری مثلثی}
\itm
یک پیمایش کامل
\fn{full walk}
از T رأس‌ها را در هر بار اولین ملاقات و همچنین هنگام بازگشت از هر فرزند لیست می‌کند.
پیمایش کامل این درخت را W می‌نامیم.
از آن‌جا که هر یال درخت دقیقاً دو بار طی می‌شود، داریم:
$$c(W) = 2c(T)$$
همچنین می‌دانیم که $c(T) \leq c(H^*)$ پس:
$$c(W) \leq 2c(H^*)$$
\end{itemframe-s}


\begin{itemframe-s}{مسئله فروشنده دوره گرد}{حل با نابرابری مثلثی}
\itm
حال می‌خواهیم از روی W یک دور همیلتونی بسازیم. به‌دلیل نا‌برابری مثلثی، حذف یک رأس هزینه را افزایش نمی‌دهد.
بنابراین بازدیدهای تکراری را حذف می‌کنیم، به طوری که در نهایت هر رأس تنها یک بار در لیست باشد.
این همان پیمایش پیش تریتب درخت است.
\itm
دور همیلتونی‌ای که بدین ترتیب ساخته می‌شود را H می‌نامیم. اما چرا و H دور همیلتونی است؟ چون هر رأس را دقیقاً یکبار دارد و توجه کنید که گراف کامل است بنابراین هر ترتیب دلخواه از رأس‌های دور است.
\itm
همچنین دیدم که:
$$c(H) \leq c(W)$$

\end{itemframe-s}


\begin{itemframe-s}{مسئله فروشنده دوره گرد}{حل با نابرابری مثلثی}
\itm
ترکیب $c(H) \leq c(W)$ و $c(W) \leq 2c(H^*)$ می‌دهد:
$$c(H) \leq 2c(H^*)$$
که برهان را کامل می‌کند.
\itm
ضریب تقریب این الگوریتم قابل قبول است اما معمولاً بهترین انتخاب برای این مسأله نیست.
الگوریتم‌های تقریبی دیگری هم وجود دارند که در عمل بسیار بهتراند.
\itm
همچنین ثابت می‌شود اگر
$P \neq NP$
باشد، آنگاه
هیچ الگوریتم تقریبی چندجمله‌ای
برای این مسأله‌ی وجود ندارد. (فارغ از اینکه چه ضریب تقریبی مدنظر داشته باشید.)
\end{itemframe-s}